\documentclass[paper=a4, fontsize=10pt]{jlreq}

\usepackage{amsmath, amssymb}
\usepackage{enumerate}
\usepackage{tikz}
\usepackage{listings, xcolor}

\lstset{
basicstyle = {\ttfamily},
frame = {tbrl},
breaklines = true,
numbers = left,
showspaces = false,
showstringspaces = false,
showtabs = false,
keywordstyle = \color{blue},
commentstyle = {\color[HTML]{1AB91A}},
identifierstyle = \color{black},
stringstyle = \color{brown},
captionpos = t
}

\title{経済学基礎 課題2}
\author{学籍番号 1280391 \\ 細川 夏風}
\date{\today}

\begin{document}

\maketitle

\section*{問題 2-1}
\begin{enumerate}[1) ]
  \item 約75％（74.9％）高くなる。\\
  大卒は高卒よりも教育年数が4年多く、$1.15^4 \approx 1.749$ となるため。

  \item 卒業後9年目（9年後）。\\
  大卒所得は高卒より46\%高い（$1.10^4 = 1.46$）。高卒所得を$x$とし、$y$年目の生涯所得差が0を上回る条件は $(1.46x - x)y - 4x > 0 \Rightarrow 0.46y > 4 \Rightarrow y > 8.69$ となるため。
\end{enumerate}

\section*{問題 2-2}
自然実験。\\
人の命に関わる法律の適用を、無作為抽出によって人為的にグループ分けすることは倫理的にも現実的にも不可能だからである。

\section*{問題 2-3}
不可能。\\
平均値はすべてのデータを足してデータの数で割ったものであり、全員が平均点を上回ることは数学的にあり得ない。

\section*{問題 2-4}
\begin{enumerate}[1) ]
  \item アパート3\\
  総費用（家賃＋通勤時間×機会費用）を計算すると、アパート3（59,500円）が最も安いため。

  \item 1時間当たり2,500円より高い（$x > 2500$）\\
  機会費用を$x$円/時間とすると、$62,000 + 1x < 52,000 + 5x \Rightarrow 10,000 < 4x \Rightarrow x > 2500$ となるため。
\end{enumerate}

\section*{問題 2-5}
同意しない。\\
経済学では純便益（便益－費用）が最大になるポイントで行動を止めるのが合理的と考える。全力を尽くすことは、限界費用が限界便益を上回ってもリソースを投入し続けることを意味するため。

\section*{問題 2-6}
中心部の家賃に関する直接費用を下げるか、郊外からの通勤における間接費用（機会費用）を上げる。\\
具体例：中心部居住者への家賃補助、郊外からの通勤に対するガソリン税引き上げや渋滞税導入など。

\section*{問題 2-7}
最適な座席: C席、最適な駐車場: M駐車場\\
限界分析により、座席はCからBに変更すると限界費用(6ドル)が限界便益(5ドル)を上回る。駐車場はMからNに変更すると限界費用(20ドル)が限界便益(10ドル)を上回るため、それぞれ状況が悪化する手前が最適となる。

\end{document}